\documentclass[letterpaper,11pt]{report}
\usepackage[utf8]{inputenc}
\usepackage[default]{raleway}

%\usepackage{fontspec}
\usepackage{fontawesome}

\usepackage{amsmath}

\usepackage[letterspace=200]{microtype}

\usepackage{multicol}

\usepackage{anyfontsize}

\usepackage{float}

\usepackage{supertabular}
%\usepackage{longtable}
%\usepackage{tabu}

\usepackage{array}

\usepackage[table]{xcolor}
\newcounter{rowcount}
\setcounter{rowcount}{0}

\usepackage{tikz}
\usetikzlibrary{fadings}

\usepackage{geometry}

\geometry{
	right=15mm,
	left=15mm,
	top=25mm,
	bottom=20mm
}

\usepackage{titlesec}
\titlespacing*{\section}{0cm}{0.8cm}{0.15cm}
\titleformat*{\section}{\lsstyle \color[rgb]{0.2,0,0} \bf \fontsize{12}{12}\selectfont}

\usepackage{cellspace}
%\cellspacetoplimit 30cm
%\cellspacebottomlimit 30pt
\setlength{\tabcolsep}{10pt}

\definecolor{cafe}{rgb}{0.250980392,0,0}
\definecolor{verde}{rgb}{0,0.5,0}
\definecolor{azul}{rgb}{0,0,0.5}

\setlength{\columnsep}{1cm}
\columnseprule 1pt
\def\columnseprulecolor{\color{gray!50}}

\pagestyle{empty}


\newcommand{\linea}[1]{
	
	\begin{tikzpicture}[remember picture, overlay] 

	\draw[line width=1pt, color=cafe] (0,0.4) -- (0,-#1);

	\end{tikzpicture}
}

\usepackage{hyperref}
\hypersetup{
	colorlinks=false,
	hidelinks=true,
}

% Evitar espacios grandes entre títulos de secciones y párrafos.
\setlength{\parskip}{0pt}

% Elementos con línea vertical y bullet.
\newlength{\sepVertical}
\setlength{\sepVertical}{21pt}
\newcommand{\elementoInicio}[2]{\noindent\hfill\parbox[t]{\columnwidth}{\parbox[t]{0.12\columnwidth}{\raggedleft #1}\hfill{\color{cafe}\vrule width1pt height 10pt}\llap{\Large$\color{cafe}\bullet \hspace{-3.1pt}$}\hfill\parbox[t]{0.8\columnwidth}{#2}}%

\vspace{-1.1pt}
}
\newcommand{\elementoMitad}[2]{\noindent\hfill\parbox[t]{\columnwidth}{\parbox[t]{0.12\columnwidth}{\raggedleft #1}\hfill{\color{cafe}\vrule width1pt height \sepVertical}\llap{\Large$\color{cafe}\bullet \hspace{-3.1pt}$}\hfill\parbox[t]{0.8\columnwidth}{#2}}%

\vspace{-1.1pt}
}
\newcommand{\elementoFin}[2]{\noindent\hfill\parbox[t]{\columnwidth}{\parbox[t]{0.12\columnwidth}{\raggedleft #1}\hfill{\color{cafe}\rule[-3pt]{1pt}{\sepVertical+3pt}}\llap{\Large$\color{cafe}\bullet \hspace{-3.1pt}$}\hfill\parbox[t]{0.8\columnwidth}{#2}}}




\begin{document}




\begin{tikzpicture}[remember picture, overlay] 

	\useasboundingbox (0cm,0cm) rectangle (\paperwidth,5cm);

	\node[anchor=north west, minimum width=\paperwidth, minimum height=5cm, fill=brown!340, inner sep=0pt, align=center] at (current page.north west) {
		\\[40pt] \textcolor{white}{\fontsize{24}{24}\selectfont {\bf JOSÉ ALBERTO} LÓPEZ LÓPEZ} \\[10pt] 
		\textcolor{white}{\large BACHELORS IN PHYSICS} \\[5pt] 
		\textcolor{white}{\large MASTERS IN ARTIFICIAL INTELLIGENCE} \\[5pt] 
		\textcolor{gray}{\href{mailto:JLL.6@hotmail.com}{\faEnvelope\ \ JLL.6@hotmail.com}}
	};

\end{tikzpicture}


\vspace{2cm}


%\begin{multicols}{1}

%\begin{table}[H]
%\begin{tabular}{b{0.25\textwidth} b{0.75\textwidth}}

%\includegraphics[scale=0.2]{J} & 

\section*{ABOUT ME \vskip 3pt \hrule width0.6\textwidth} 
	
	I belive that physics and AI are two careers which complement each other very well and both are powerful tools for solving a great number of problems. I consider myself a dedicated, responsable, and very capable person, with the ability to quickly learn new things. I like to play sports and I have great passion for music.
	
\vspace{0.35cm}

%\end{tabular}
%\end{table}

%\end{multicols}

\begin{multicols*}{2}

\section*{EXPERIENCE \vskip 3pt \hrule width0.6\columnwidth}%

\setlength{\tabcolsep}{0pt}

\begin{tabular*}{\columnwidth}{>{\raggedleft}p{0.085\columnwidth} >{\centering \Large}p{0.08\columnwidth} p{0.83\columnwidth}}
	2018 & $\color{cafe} \bullet$\llap{\linea{4.16}} & {\bf Algorithms and Prototipes Developer} \newline\textsc{Artificial Intelligence Research Center} $\cdot$ University of Veracruz \ \href{https://maps.google.com/?cid=17992863989216810333}{\textcolor{red}{\faMapMarker}} \newline \color{black!70} C++ developer of various company projects, involving computer vision, autonomous movement detection in vehicles and embedded systems using accelerometers and Furier analisys. \\[-0.3cm] \tabularnewline

	2017 & $\color{cafe} \bullet$ & {\bf Engineer Assistant} \newline\textsc{Mechanics and Reactor Engineery departments} $\cdot$ Laguna Verde Nuclear Power Plant \ \href{https://maps.app.goo.gl/aeNEVSTLKJTFyx4E9}{\textcolor{red}{\faMapMarker}} \newline \color{black!70} Engineer assistant in the 18$^{\text{th}}$ fuel recharge of Laguna Verde Nuclear Power Plant. Assisting machanical engineers in maintainance of diverses reactor machines and systems, and doing administrative work of the department. As well as performing activities in reactor engineery department like reactor monitoring, data analysis of reactor jet pumps stability among others.
\end{tabular*}

%\columnbreak
\vspace{-0.2cm}

\section*{EDUCATION \vskip 3pt \hrule width0.6\columnwidth}

\setlength{\tabcolsep}{10pt}
\begin{tabular*}{\columnwidth}{p{0.08\columnwidth} S{p{0.8\columnwidth}}}
    2020 & {\bf Masters in artificial intelligence} \newline\textsc{Artificial Intelligence Research Center} $\cdot$ University of Veracruz, Mexico \ \color{cyan!65} \href{https://www.uv.mx/dia/}{\faUniversity} \\[-0.2cm] \tabularnewline
    
    2018 & {\bf Bachelors in physics} \newline\textsc{Faculty of Physics} $\cdot$ University of Veracruz, Mexico \ \color{cyan!65} \href{https://www.uv.mx/ffia/}{\faUniversity} \newline \color{black!70} Best marks of generation. \\ 
\end{tabular*}


\section*{PROFILE \vskip 3pt \hrule width0.6\columnwidth}


My background in nature and computer sciences areas has allowed me to develop good skills to solve problems of many kinds, especially those that involve programing and mathematics. 
I have an outstanding analytical skill and I always focus on offering solutions in the most efficient way possible, but also looking for the simplest, tasteful and time-proof results. \\[-0.6cm]

\section*{PROGRAMMING LANGUAGES \newline \& TECHNOLOGIES \vskip 6pt \hrule width\columnwidth}

{\renewcommand{\arraystretch}{1.56}
\def\barheight{6pt}
\vspace{-0.25cm}
\begin{tabular}{>{\bf} m{0.25\columnwidth} m{0.568\columnwidth}}

	 & \color{gray} \bf Junior \hfill \color{cafe} Senior \\

    C++ & \begin{tikzpicture}[line width=5pt] 
    \shade[left color=gray!60, right color=gray!60] (0,0) rectangle (5cm,\barheight); 
    \shade[left color=brown!300, right color=gray!60] (3.95cm,0) rectangle (4.7cm,\barheight);
    \shade[left color=brown!300, right color=brown!300] (0,0) rectangle (4.0cm,\barheight); 
    \end{tikzpicture} \\
    
    Python  & \begin{tikzpicture}[line width=5pt] 
    \shade[left color=gray!60, right color=gray!60] (0,0) rectangle (5cm,\barheight); 
    \shade[left color=brown!300, right color=gray!60] (3.45cm,0) rectangle (4.2cm,\barheight);
    \shade[left color=brown!300, right color=brown!300] (0,0) rectangle (3.5cm,\barheight); 
    \end{tikzpicture} \newline {\bf \small \color{gray} numpy, pandas, statsmodels, keras, sklearn, tkinter, etc.} \\
    
    Mathematica & \begin{tikzpicture}[line width=5pt] 
    \shade[left color=gray!60, right color=gray!60] (0,0) rectangle (5cm,\barheight); 
    \shade[left color=brown!300, right color=gray!60] (3.45cm,0) rectangle (4.2cm,\barheight);
    \shade[left color=brown!300, right color=brown!300] (0,0) rectangle (3.5cm,\barheight); 
    \end{tikzpicture} \\
    
    SQL, MySQL & \begin{tikzpicture}[line width=5pt] 
    \shade[left color=gray!60, right color=gray!60] (0,0) rectangle (5cm,\barheight); 
    \shade[left color=brown!300, right color=gray!60] (3.05cm,0) rectangle (3.8cm,\barheight);
    \shade[left color=brown!300, right color=brown!300] (0,0) rectangle (3.1cm,\barheight); 
    \end{tikzpicture} \\
    
    Git & \begin{tikzpicture}[line width=5pt] 
    \shade[left color=gray!60, right color=gray!60] (0,0) rectangle (5cm,\barheight); 
    \shade[left color=brown!300, right color=gray!60] (3.05cm,0) rectangle (3.8cm,\barheight);
    \shade[left color=brown!300, right color=brown!300] (0,0) rectangle (3.1cm,\barheight); 
    \end{tikzpicture} \\
    
    Java & \begin{tikzpicture}[line width=5pt] 
    \shade[left color=gray!60, right color=gray!60] (0,0) rectangle (5cm,\barheight); 
    \shade[left color=brown!300, right color=gray!60] (2.45cm,0) rectangle (3.2cm,\barheight);
    \shade[left color=brown!300, right color=brown!300] (0,0) rectangle (2.5cm,\barheight); 
    \end{tikzpicture} \\
    
    Fortran & \begin{tikzpicture}[line width=5pt] 
    \shade[left color=gray!60, right color=gray!60] (0,0) rectangle (5cm,\barheight); 
    \shade[left color=brown!300, right color=gray!60] (2.45cm,0) rectangle (3.2cm,\barheight);
    \shade[left color=brown!300, right color=brown!300] (0,0) rectangle (2.5cm,\barheight); 
    \end{tikzpicture} \\
    
    Javascript & \begin{tikzpicture}[line width=5pt] 
    \shade[left color=gray!60, right color=gray!60] (0,0) rectangle (5cm,\barheight); 
    \shade[left color=brown!300, right color=gray!60] (1.95cm,0) rectangle (2.7cm,\barheight);
    \shade[left color=brown!300, right color=brown!300] (0,0) rectangle (2cm,\barheight); 
    \end{tikzpicture} \\
    
    HTML, CSS & \begin{tikzpicture}[line width=5pt] 
    \shade[left color=gray!60, right color=gray!60] (0,0) rectangle (5cm,\barheight); 
    \shade[left color=brown!300, right color=gray!60] (1.95cm,0) rectangle (2.7cm,\barheight);
    \shade[left color=brown!300, right color=brown!300] (0,0) rectangle (2cm,\barheight); 
    \end{tikzpicture} \\
    
    \LaTeX & \begin{tikzpicture}[line width=5pt] 
    \shade[left color=gray!60, right color=gray!60] (0,0) rectangle (5cm,\barheight); 
    \shade[left color=brown!300, right color=gray!60] (3.95cm,0) rectangle (4.7cm,\barheight);
    \shade[left color=brown!300, right color=brown!300] (0,0) rectangle (4.0cm,\barheight); 
    \end{tikzpicture} \\
    
    Others & \begin{tikzpicture}[line width=5pt] 
    \shade[left color=gray!60, right color=gray!60] (0,0) rectangle (5cm,\barheight); 
    \shade[left color=brown!300, right color=gray!60] (1.95cm,0) rectangle (2.7cm,\barheight);
    \shade[left color=brown!300, right color=brown!300] (0,0) rectangle (2cm,\barheight); 
    \end{tikzpicture} \\

\end{tabular}
}

\columnbreak

\section*{PROJECTS \vskip 3pt \hrule width0.6\columnwidth}

{\bf Robot Hostess} \ {\color{red} \href{https://www.youtube.com/watch?v=KU4NL_bc_rA}{\faYoutubePlay}} \ {\href{https://github.com/JoseAlbertoLopezLopez/CV_projects/tree/master/RobotHostess}{\faGithub}} \\
{\color{black!70} I programmed the entire code of a Robot hostess in C++ with ROS. The robot detects diners via face recognition and lead them to a suitable table.}\\

\noindent {\bf Augmented reality} \ {\color{red} \href{https://www.youtube.com/watch?v=aLcEEVrU13I}{\faYoutubePlay}} \ {\href{https://github.com/JoseAlbertoLopezLopez/CV_projects/tree/master/AugmentedReality}{\faGithub}} \\
{\color{black!70} Example of augmented reality made with Python and OpenCV.}\\

\noindent {\bf Traveling salesman problem visualization} \ {\color{red} \href{https://www.youtube.com/watch?v=VK9NzBgPXBw}{\faYoutubePlay}} \ {\href{https://github.com/JoseAlbertoLopezLopez/CV_projects/tree/master/SimulatedAnnealing}{\faGithub}} \\
{\color{black!70} Visualization of traveling salesman problem solved with simulated annealing.}\\

\noindent {\bf Grasp Objects in Clutter (Masters Thesis)} \\
{\color{black!70} I developed an algorithm for accommodate objects in order to facilitating later grasps by defining grasping and neighborhood rules, and appling AI nature inspired algorithms.}\\

\noindent {\bf Automatic Car's Anomalous Movements Detection (Bachelors Thesis)} \ {\href{https://github.com/JoseAlbertoLopezLopez/CV_projects/tree/master/BachelorsThesis}{\faGithub}} \\
{\color{black!70} I made a program for detecting several kinds of movements in cars using acceleration sensors and signal procesing theory whose data were processed by a neutral network.}

\vspace{0.5cm}

\section*{\bf CERTIFICATIONS \vskip 3pt \hrule width0.6\columnwidth}

I usually don't take courses, I read books about the subjects I need.\\

{\setlength{\tabcolsep}{0pt}%
\renewcommand{\arraystretch}{1.9}%
\rowcolors{0}{}{gray!20}%
\par\centering%
\begin{tabular}{>{\centering \bf} m{0.31\columnwidth} >{\centering \bf} m{0.31\columnwidth} >{\centering \bf} m{0.31\columnwidth}}%
%\hline
\rowcolor{cafe}{\color{white} Institution} \vspace{0.06cm} & {\color{white} Subject} \vspace{0.06cm} & {\color{white} ID and URL} \vspace{0.06cm} \tabularnewline
%\hline
Udemy & Time Series Data Analisis & \href{http://ude.my/UC-51b1bb6c-b324-41d7-b11f-3826f7999574}{\Large\color{gray}\faShareSquareO} \tabularnewline
%\hline
Udemy & Git & \href{http://ude.my/UC-c48e4dd1-59c9-4f2c-b199-c9d09a74f9e8}{\Large\color{gray}\faShareSquareO} \tabularnewline
%\hline
Udemy & SQL: Basic to Advanced & \href{http://ude.my/UC-6acc8434-7c34-40b5-ae1b-b813b2252e33}{\Large\color{gray}\faShareSquareO} \tabularnewline
%\hline
SoloLearn & SQL & \href{https://www.sololearn.com/Certificate/1060-849693/pdf/}{\Large\color{gray}\faShareSquareO} \tabularnewline
%\hline
SoloLearn & Java & \href{https://www.sololearn.com/Certificate/1060-849693/pdf/}{\Large\color{gray}\faShareSquareO} \tabularnewline
%\hline
\end{tabular}\par}

\columnbreak

\section*{\raggedright \bf COMPLEMENTARY\ \ EXPERIENCE \vskip 3pt \hrule width\columnwidth}

% Espacio entre elementos original: \\[0.35cm].
\setlength{\sepVertical}{26.6pt}
\elementoInicio{\footnotesize 2 years}{Python libraries for data science like NumPy, Pandas, ScikitLearn, Statsmodels, FacebookProphet, among others.}

\elementoMitad{\footnotesize 2 years}{Mathematical foundations of models for time series analisys, like ARIMA, SARIMA, VARMA, Dickey-Fuller test, Granger causality, Holt-Winters method, EWMA, ETS, Hodrick-Prescott filter, etc.}

\elementoMitad{\footnotesize 2 years}{Experience in parallel programing with OpenMS and MSI libraries in C++, C and Fortran.}
\elementoMitad{\footnotesize 2 years}{Computer vision with OpenCV library in C++ and Python.}
\elementoMitad{\footnotesize 2 years}{Robotic Operative System (ROS) with C++ and Python.}
\elementoMitad{\footnotesize 2 years}{Experience in Keras and TensorFlow for deep learning and machine learning.}
\elementoMitad{\footnotesize 3 years}{Mathematical foundations of dense, convolutional, and recurrent neural networks; as well as evolutionary and nature-inspired algorithms.}
\elementoMitad{\footnotesize 5 years}{Large experience in data visualization with Gnuplot, Tikz-PGF library and Matplotlib (Python).}
\elementoMitad{\footnotesize 7 years}{Use of Linux systems (terminal, etc.).}

\elementoMitad{2017}{Engineery of Nuclear Pants Course, at Laguna Verde Nuclear Power Plant.}
\elementoMitad{2017}{Nuclear Security Course, at Laguna Verde Nuclear Power Plant.}

\elementoMitad{2016}{VII Physics and Nano-structures Workshop, organized by Nanosciences and Nanotechnology Center of UNAM, Ensenada, B.C, Mexico. \newline \hfill\color{black!70} Working in a quantum computing theory subject.}

\setlength{\sepVertical}{21pt}
\elementoInicio{2015}{Advanced Summer School 2015, organized by Research and Advenced Studies Center of IPN (National Polytechnic Institute), Cinvestav-IPN.
\newline\color{black!70} Advanced topics of math, physics and engineering.}

\elementoFin{2015}{XXV Summer of Science, organized by Mexican Academy of Sciences, at the Institute of Nuclear Sciences of the UNAM (Autonomous National University of Mexico). \newline \color{black!70} Working with data science topics applied to cities databases.}

%\columnbreak

\section*{\raggedright OTHER LANGUAGES \vskip 3pt \hrule width0.8\columnwidth}


\begin{tabular*}{\columnwidth}{>{\raggedright} p{0.185\columnwidth} >{\centering} p{0.05\columnwidth}}
    {\bf Spanish} & (native)
\end{tabular*}


\section*{CONTACT \vskip 3pt \hrule width0.6\columnwidth}


\begin{itemize}
	\item[\fontsize{11}{11}\selectfont \href{mailto:JLL.6@hotmail.com}{\faEnvelopeO}] \href{mailto:JLL.6@hotmail.com}{JLL.6@hotmail.com}%
	\item[\fontsize{13}{13}\selectfont \faMobile \hspace{0.07cm}] +52 2281054832%
	\item[\fontsize{12}{12}\selectfont \color{azul} \faPhone] +52 2969647471%
	\item[\fontsize{12}{12}\selectfont \color{azul} \href{https://www.linkedin.com/in/alberto-l\%C3\%B3pez-177a281a0/?locale=en_US}{\faLinkedinSquare}] \href{https://www.linkedin.com/in/alberto-l\%C3\%B3pez-177a281a0/?locale=en_US}{Alberto López}%
\end{itemize}

%\vfill\mbox{}
%\vspace{0.25cm}
%\columnseprule 0pt

\section*{REFERENCES \vskip 3pt \hrule width0.6\columnwidth}

\noindent {\bf Marxlenin Zapata Yañez} \newline {\color{black!70} Chief of Reactor Engineering department at Laguna Verde Nuclear Power Plant. \newline
Relation: Previous boss.} \newline \href{mailto:marxlenin.zapata@cfe.gob.mx}{\faEnvelopeO\ marxlenin.zapata@cfe.gob.mx}%

\vspace{0.3cm}%

\noindent {\bf Irving Morales Agiss, PhD} \newline {\color{black!70} Director of data science MORLANmx enterprise. \newline
Relation: Professor.} \newline \href{mailto:irvingfisica@gmail.com}{\faEnvelopeO\ irvingfisica@gmail.com}%

\vspace{0.3cm}%

\noindent {\bf Antonio Marín Hernández, PhD} \newline {\color{black!70} Researcher at the Artificial Intelligence Research Center of University of Veracruz, Mexico. \newline
Relation: Thesis assesor.} \newline \href{mailto:anmarin@uv.mx}{\faEnvelopeO\ anmarin@uv.mx}

\vspace{0.3cm}

%\noindent {\bf Sergio Lerma Hernández, PhD} \newline {\color{black!70} Researcher at Faculty of Physics of University of Veracruz, Mexico. \newline
%Relation: Professor.} \newline \href{mailto:sergioalerma@gmail.com}{\faEnvelopeO\ sergioalerma@gmail.com}

\end{multicols*}

\end{document}